\documentclass[../master.tex]{subfiles}

\begin{document}
\chapter*{はじめに}
物理の面白さの1つは一見すると難しそうに見える問題が見方を変えるだけで驚くほど容易に解けるようになる点にある。
例えば、相互作用のある力学系の最小の例として質点2つの衝突を考えると、
エネルギー保存則と運動量保存則から連立方程式を立てて直接計算することもできるが、
計算は煩雑になりがちである。
ところが重心座標と相対座標へと変数を取り替えると、
この問題は互いに独立な一体問題へと還元され、
あとはシンプルな公式を適用するだけで運動が求まる。
このように系を表す自由度を適切に選び直すことによって、
系のダイナミクスは見通しよく記述できる。
しかし多体系においては「何が適切な見方なのか」自体が容易ではない。

多体系では系を構成する各自由度の運動を手計算で追跡することは現実的でなく、
仮に行うとしても分子動力学シミュレーションのような数値計算に頼らざるを得ない。
数値計算によって得られる軌道を可視化することは直観を与える一方で、
それだけで系の物理を理解したと言うのは難しい。
理解に向けては、得られた多数の自由度をそのまま保持するのではなく、
位置や運動量の統計量など、より少数の量に情報を整理して観察を進める必要がある。
すると例えば、壁近傍で速度が低下する傾向や、
ゆらぎがある1つの尺度で整理されるといった、
個々の軌道の詳細を越えた規則性が現れてくる。

実際、この種の規則性は大規模計算が不可能であった19世紀にも流体力学や熱力学として体系化されてきた。
両者に共通しているのは微視的な質点の描像に固執するのではなく、
粗視化によって導入される自由度(場や少数のマクロ変数)へと記述を切り替えることで、
多体系の振る舞いを見通しよく捉えるという態度である。
粗視化により系を理解するのにふさわしい自由度が得られると同時に、
その自由度で記述された状態がいくつかの「型」に分かれることにも気づく。
この状態の型分けが相であり、多体系を理解するとは、
微視的模型から出発して適切な自由度による描像へと写し替え、
最終的に相の分類とその性質・起源を明らかにすることだと言える。

\bigskip
相を理解する枠組みは、20世紀を通して段階的に拡張されてきた。
その古典的な出発点は1930年代のLandau理論であり、
相を自発的対称性の破れと局所秩序変数によって特徴づける見方が確立された。
この枠組みは物性物理における相転移の理解の基礎を与えただけでなく、
1960年代のHiggs機構を通して素粒子物理における真空の理解にも深く関わっている。
物性では相を、素粒子では真空を扱っているという違いはあるが、
いずれの場合も「対称性がどのように実現されているか」によって状態を分類するという視点が中心にあった。
したがって、現代の量子多体系をめぐる多くの議論も、まずはこのLandau的な視点から出発している。

しかし1960年代末以降、量子論的・位相的な効果を正面から扱う必要が生じるにつれて、
Landau理論だけでは見通せない現象が次第に前景化した。
場の理論では、1969年の軸性アノマリーの発見以降、
古典的には保存されるはずの対称性が量子論では破れるという事実が、
対称性の理解そのものを拡張する重要な契機となった。
対称性は単にハミルトニアンや作用の不変性を述べるだけではなく、
許される低エネルギー有効理論や境界自由度、さらには実現可能な真空構造そのものを拘束する量として理解されるようになった。
一方、物性物理では1980年代の整数量子ホール効果と分数量子ホール効果の研究を通して、
局所秩序変数を用いずとも安定に区別される量子相が現実に存在することが明確になった。
その後、2010年前後にはSPT相の枠組みも整備され、
対称性の破れとは異なる形で対称性が量子相を特徴づける仕方が体系的に理解されるようになった。

この2つの流れを媒介したのが、1980年代末から1990年代にかけて整備されたトポロジカル場の理論である。
Chern-Simons理論は量子ホール系の有効理論やアノマリー流入の記述において中心的な役割を果たし、
BF理論はHiggs相や離散ゲージ理論の低エネルギー記述として現れる。
この時期に発展したTQFTの考え方は、
局所的な秩序変数ではなく、大域的なトポロジーや境界、応答係数によって相を記述する見方を洗練した。
ここで重要なのは、物性と素粒子で用いられる対象が異なっていても、
背後にある数学的言語がかなりの程度共通しているという点である。
私がこの周辺を面白いと思う理由の1つは、別々の問題意識から出発したはずの2つの分野が、
気づくと同じ理論を異なる言葉で使っているところにある。

こうしたトポロジカル場の理論が与えた最も印象的な描像の1つが、
トポロジカルな励起、すなわちエニオンの描像である。
2次元では粒子交換の統計性がボソンとフェルミオンに限られず、
より一般の位相やユニタリ変換として実現されうる。
量子ホール系の研究を通してエニオンは具体的な物理対象として前景化し、
トポロジカル相の性質を励起の統計性から理解する視点が生まれた。
この視点では相の違いは単に秩序変数の有無ではなく、
どのような準粒子が現れ、それらがどう交換(braiding)・融合(fusion)され、どのように動けるかによっても特徴づけられる。
現在ではこの考え方はかなり広く受け入れられているが、
それはLandau理論の否定というより、相を記述する言葉が
対称性の破れからトポロジカル励起へと拡張されてきた結果だと見る方が自然である。

他方で1980年代には、量子計算というまったく別の動機から、
量子力学的な多体性そのものに新しい意味が与えられた。
Feynmanの量子シミュレーションの提案やDeutschの普遍量子計算の議論を通して、
量子系を単に自然現象として記述する対象ではなく、
情報処理を担う装置として捉える視点が立ち上がった。
1990年代には量子アルゴリズムが発展し、
同時にShor符号、CSS符号、stabilizer formalism などの仕事によって
量子誤り訂正の基本概念が整備された
\cite{Shor1995-ii,Calderbank1996-si,Steane1996-is,Gottesman1997-ep}。
ここで重要なのは量子誤り訂正が単なる工学的補助手段ではなく、
多体系の量子状態空間の中に情報をどのように埋め込み、
局所的な擾乱からどのように守るかという、きわめて物理的な問題として立ち現れたことである。

この流れの中で2000年代に現れたKitaevのtoric codeは、
量子情報と量子多体系の接点を決定的に可視化した模型だったと言ってよい。
toric codeは量子誤り訂正符号であると同時に、
局所可換ハミルトニアンの基底状態として実現されるトポロジカル秩序相でもあり、
その励起はエニオンとして理解される。
すなわちこの模型では、基底状態の位相的縮退、局所演算子に対する頑健性、
準粒子の生成と移動、シンドローム測定と誤り訂正といった概念が1つの言語で結びついている。
量子情報の側から見れば情報保護の模型であり、
量子多体系の側から見ればトポロジカル相の可解模型であるという二重の読みが可能になったことが、
その後の発展に大きな影響を与えた。

2000年代から2010年代にかけて量子情報の文脈では、
このtoric codeを出発点として局所可換ハミルトニアンを系統的に構成する方法論が発展した。
表面符号やその高次元化、安定化符号の一般論、hypergraph product符号のような構成法、
そして2次元局所可換符号では自己修復量子メモリが難しいことを示すno-go定理の整理などを通して、
量子情報は意図せずして可解な量子スピン模型の豊かな構成論を手にしたのである。
この構成論の面白さは、情報保護の性能を追うことがそのまま、
どのような励起構造を持つ局所模型が可能かを問うことにつながっている点にある。
とくにno-go定理を回避するために3次元以上へ目を向けると、
2次元のエニオン像では記述しきれない挙動を示す模型が自然に現れる。

その代表が2011年のHaah符号であり、
そこで現れた点状励起は単独では自由に動けず、
従来のひも演算子によるエニオンの描像から大きく外れていた。
この種の励起は、2011年から2016年頃にかけて次第に整理され、
fracton, lineon, planonといった、低次元的な運動性しか持たない励起として理解されるようになった。
この時点で重要なのは、トポロジカルな励起の概念が2次元エニオンに尽きないこと、
そして励起の「統計性」だけでなく「可動性」そのものが相を特徴づける本質的なデータであることが明瞭になったことである。
したがって、現在の量子多体系では相を対称性と励起の両面から見るだけでなく、
励起がどの方向にどの演算子によって、どの程度局所的に動けるかという運動学的構造まで含めて考える必要がある。

他方で場の理論の側でも、2010年代に入ると一般化対称性の言語が急速に整備された。
これは、通常の対称性(0-form対称性)だけでは捉えきれない保存則や欠陥演算子を、
より高次の演算子に作用する対称性として捉え直す枠組みである。
この見方はアノマリーやTQFTの理解を整理し直すだけでなく、
トポロジカル相や拘束された励起を対称性の立場から統一的に記述する道を開いた。
さらに2010年代後半以降、fracton相や多極子保存則を持つ系が
高階テンソルゲージ理論や多極子代数の言葉で記述されるようになると、
場の理論で育ってきた一般化対称性の発想と、
格子模型で見出されていた特異な励起構造とが次第に合流し始めた。
% 私の理解では、現在はまだこの翻訳作業の途上にあり、
% 物性の人が励起や模型の言葉で眺めてきたものと、
% 素粒子の人が対称性やアノマリーの言葉で眺めてきたものとが、
% ようやく同じ図の上に描かれつつある段階にある。

この周辺の話題はしばしば「量子情報」という大きな枠組みのもとでひとまとめに語られるが、
その内部には関心の異なる複数の流れがある。
量子計算の文脈では、量子系をいかに制御し、情報処理や誤り耐性に役立てるかが主たる関心となる。
これに対して量子多体系の文脈では、量子状態がどのような相を実現し、
どのような励起やエンタングルメント構造を持つかを理解することが主題であり、
計算資源としての有用性は必ずしも第一の関心ではない。
両者は局所性、エンタングルメント、トポロジカル相、量子誤り訂正符号といった概念を共有しており密接に関わっているが、
何を面白いとみなすかは同一ではない。
そして、近年は量子情報の概念を量子多体系へ持ち込む流れはかなり見えやすい一方で、
量子多体系の側で現れた新しい相や励起の理解を、
量子計算や量子誤り訂正の具体的な構成法へ持ち帰る方向の研究は
まだ見えにくいように思われる。
私はこの点を少しもったいないと感じており、
量子誤り訂正符号を量子計算の道具としてだけでなく、
現代的な量子多体系の描像を具体的に実現するための自然な言語として読み直したいと考えている。

\bigskip
そうした多体系の中で近年、多極子保存則を満たす系が注目されている。
多極子保存則は励起の運動を強く拘束し、
局所秩序変数や通常の対称性(より正確には局所演算子に作用する0-form対称性)だけでは捉えにくい構造をもつ量子相を生みうるため、
物性物理および場の理論の双方から研究が進められている。

多極子保存則を満たす系についての研究は
数値計算によるもの
\cite{Sala2020-tm,Khemani2020-uw,Moudgalya2022-pg,Lake2022-nn}
と、場の理論によるもの
\cite{Pretko2017-jy, Gromov2019-mm,
    Zhai2021-nt, Grosvenor2021-zl, Ebisu2025-vv}
が行われている。
数値計算による研究では、熱化やエルゴード性と併せて具体的な模型の性質が調べられてきた。
しかし、系が可解ではないため得られる基底状態やそのダイナミクスの解釈が難しく、
議論が個別模型ごとの解析にとどまりやすい傾向がある。

場の理論を用いた研究では新しいクラスの対称性の一種として捉え、
作用からスタートして対称性と系のトポロジカルな性質を調べることを中心としている。
しかしこれは抽象度が高い理論になっているため、
その場の理論で記述するのがふさわしい系で生じる多体現象の説明には適している一方、
それに対応する具体的な微視的格子模型を提示することは必ずしも容易でない。
この場の理論を用いた研究は特定の対称性を持った系の持つ相の分類を目的としているが、
相の分類を具体的な模型に結びつけて検証しようとすると、
解析可能な例が限られていること自体がボトルネックになっているように思われる。

そこで本研究では多極子保存則を満たす量子多体系について、
量子誤り訂正符号の分野で発展した構成法のうち、
hypergraph product符号の方法を利用し、
解析可能な格子模型を構成した。
この模型を準粒子とゲージ場の言葉へ写像することで、
トポロジカルな相の構造を見通しよく読み解いた。
これにより多極子保存則を持つ系について、相互作用を明示した微視的多体模型から、
系の振る舞いをよく説明する適切な見方での相まで追跡して理解する道筋が得られる。
この模型ではトポロジカルな基底状態の縮重度や準粒子の可動性を明示的に示すことができ、
同様の手法により具体的な模型を系統的に構成するための見通しも得られる。
以上により、多極子保存則を満たす系における相の理解と分類に向けた足場を与える。

\bigskip


\end{document}
